\documentclass[12pt,a4paper]{article}

\usepackage[utf8]{inputenc}
\usepackage[T1]{fontenc}
\usepackage[spanish]{babel}
\usepackage{geometry}
\usepackage{booktabs}
\usepackage{longtable}
\usepackage{listings}
\usepackage{xcolor}
\usepackage{hyperref}
\usepackage{graphicx}
\usepackage{array}
\usepackage{tabularx}
\usepackage{titlesec}
\usepackage{fancyhdr}
\usepackage{parskip}

\geometry{
  top=2.5cm,
  bottom=2.5cm,
  left=3cm,
  right=2.5cm
}

\definecolor{codebg}{rgb}{0.95,0.95,0.95}
\definecolor{codeframe}{rgb}{0.7,0.7,0.7}
\definecolor{keywordcolor}{rgb}{0.0,0.3,0.6}
\definecolor{commentcolor}{rgb}{0.3,0.5,0.3}
\definecolor{stringcolor}{rgb}{0.6,0.1,0.1}

\lstset{
  backgroundcolor=\color{codebg},
  frame=single,
  rulecolor=\color{codeframe},
  basicstyle=\ttfamily\small,
  keywordstyle=\color{keywordcolor}\bfseries,
  commentstyle=\color{commentcolor}\itshape,
  stringstyle=\color{stringcolor},
  breaklines=true,
  breakatwhitespace=true,
  showstringspaces=false,
  tabsize=2,
  captionpos=b,
  numbers=left,
  numberstyle=\tiny\color{gray},
  numbersep=5pt
}

\lstdefinelanguage{yaml}{
  keywords={true,false,null,yes,no},
  keywordstyle=\color{keywordcolor}\bfseries,
  sensitive=false,
  comment=[l]{\#},
  commentstyle=\color{commentcolor}\itshape,
  stringstyle=\color{stringcolor},
  morestring=[b]',
  morestring=[b]"
}

\pagestyle{fancy}
\fancyhf{}
\rhead{\small UAMIShop -- Entrega de Proyecto}
\lhead{\small Universidad Aut\'onoma Metropolitana}
\cfoot{\thepage}
\renewcommand{\headrulewidth}{0.4pt}

\hypersetup{
  colorlinks=true,
  linkcolor=blue!60!black,
  urlcolor=blue!60!black,
  citecolor=blue!60!black,
  pdftitle={UAMIShop - Entrega de Proyecto},
  pdfauthor={UAM},
  pdfsubject={Microservicios Spring Boot}
}

\titleformat{\section}{\large\bfseries}{}{0em}{\thesection\quad}[\titlerule]
\titleformat{\subsection}{\normalsize\bfseries}{}{0em}{\thesubsection\quad}

\begin{document}

%-----------------------------------------------------------------------
% PORTADA
%-----------------------------------------------------------------------
\begin{titlepage}
  \centering
  \vspace*{1.5cm}

  {\Large\bfseries Universidad Aut\'onoma Metropolitana\par}
  \vspace{0.3cm}
  {\large Unidad Iztapalapa\par}
  \vspace{0.3cm}
  {\large Divisi\'on de Ciencias B\'asicas e Ingenier\'ia\par}

  \vspace{1.5cm}
  \rule{\linewidth}{0.8pt}\\[0.4cm]
  {\LARGE\bfseries UAMIShop\par}
  \vspace{0.3cm}
  {\Large Sistema de E-commerce con Microservicios\par}
  \rule{\linewidth}{0.8pt}

  \vspace{1.2cm}
  {\large\itshape Entrega de Proyecto Final\par}
  \vspace{0.5cm}
  {\normalsize Temas Selectos de Ingenier\'ia de Software\par}

  \vfill

  \begin{tabular}{ll}
    \textbf{Curso:}       & Temas Selectos de Ingenier\'ia de Software \\[4pt]
    \textbf{Programa:}    & Maestr\'ia en Computaci\'on y Comunicaciones \\[4pt]
    \textbf{Fecha:}       & Abril 2026 \\[4pt]
    \textbf{Versi\'on:}   & 1.0 \\
  \end{tabular}

  \vspace{1cm}
\end{titlepage}

%-----------------------------------------------------------------------
% TABLA DE CONTENIDOS
%-----------------------------------------------------------------------
\tableofcontents
\newpage

%-----------------------------------------------------------------------
% 1. INTRODUCCION
%-----------------------------------------------------------------------
\section{Introducci\'on y Objetivos del Proyecto}

UAMIShop es un sistema de comercio electr\'onico desarrollado con una arquitectura basada
en microservicios, utilizando Spring Boot 3.1 como plataforma principal de desarrollo.
El proyecto tiene como finalidad demostrar la aplicaci\'on pr\'actica de patrones de dise\~no
modernos, principios de arquitectura distribuida y t\'ecnicas de pruebas de integraci\'on en
un contexto acad\'emico.

\subsection{Objetivos Generales}

\begin{itemize}
  \item Implementar una arquitectura de microservicios funcional que cubra el flujo completo
        de un sistema de e-commerce: cat\'alogo de productos, carrito de compras, gesti\'on
        de \'ordenes y procesamiento de ventas.
  \item Aplicar patrones de dise\~no avanzados como \textit{Domain-Driven Design} (DDD),
        arquitectura hexagonal, \textit{Circuit Breaker} y el patr\'on \textit{Outbox}.
  \item Integrar un bus de mensajer\'ia as\'incrona mediante RabbitMQ para la comunicaci\'on
        entre servicios.
  \item Proveer una interfaz de usuario funcional servida a trav\'es de Nginx.
  \item Garantizar la calidad mediante pruebas de integraci\'on automatizadas con el
        framework Karate.
\end{itemize}

\subsection{Objetivos Espec\'ificos}

\begin{itemize}
  \item Configurar un API Gateway centralizado con Spring Cloud Gateway que enrute las
        peticiones entrantes hacia los microservicios correspondientes.
  \item Implementar resiliencia en el servicio de ventas mediante \textit{Circuit Breaker}
        con Resilience4j para tolerar fallos en el servicio de cat\'alogo.
  \item Garantizar consistencia eventual en el servicio de \'ordenes mediante el patr\'on
        \textit{Outbox} con publicaci\'on de eventos a RabbitMQ.
  \item Contenerizar la totalidad de la soluci\'on con Docker Compose para facilitar el
        despliegue y la reproducci\'on del entorno.
\end{itemize}

%-----------------------------------------------------------------------
% 2. ARQUITECTURA
%-----------------------------------------------------------------------
\section{Arquitectura del Sistema}

La arquitectura de UAMIShop sigue el modelo de microservicios, donde cada servicio es
aut\'onomo, posee su propia base de datos (patr\'on \textit{Database per Service}) y se
comunica con el resto a trav\'es del API Gateway (comunicaci\'on s\'incrona HTTP) y
RabbitMQ (comunicaci\'on as\'incrona por eventos).

La figura l\'ogica del sistema puede describirse de la siguiente manera:

\begin{center}
  \texttt{Cliente (Navegador)} $\rightarrow$ \texttt{Nginx / Frontend (3000)} \\
  $\downarrow$ \\
  \texttt{API Gateway (8090)} \\
  $\downarrow$ \\
  \texttt{Microservicios (8080--8083)}
  $\leftrightarrow$ \texttt{MySQL} \\
  \hspace{3.5cm} $\updownarrow$ \\
  \hspace{3.5cm} \texttt{RabbitMQ}
\end{center}

\subsection{Microservicios}

\begin{longtable}{@{}p{4.2cm}p{1.4cm}p{2.5cm}p{5.8cm}@{}}
  \toprule
  \textbf{Servicio} & \textbf{Puerto} & \textbf{Base de Datos} & \textbf{Descripci\'on} \\
  \midrule
  \endfirsthead
  \toprule
  \textbf{Servicio} & \textbf{Puerto} & \textbf{Base de Datos} & \textbf{Descripci\'on} \\
  \midrule
  \endhead
  \bottomrule
  \endfoot
  uamishop (demo)        & 8080 & uamishop\_db      & Servicio base / demo del proyecto. \\[4pt]
  uamishop-catalogo      & 8081 & catalogo\_db      & Gesti\'on de productos y categor\'ias. Expone CRUD completo con validaci\'on de SKU (\texttt{[A-Z]\{3\}-\textbackslash d\{3\}}). \\[4pt]
  uamishop-ordenes       & 8082 & ordenes\_db       & Creaci\'on y seguimiento de \'ordenes de compra. Implementa el patr\'on Outbox para publicaci\'on de eventos. \\[4pt]
  uamishop-ventas        & 8083 & ventas\_db        & Gesti\'on de carritos de compra y procesamiento de pagos. Implementa Circuit Breaker sobre el cat\'alogo. \\[4pt]
  uamishop-gateway       & 8090 & ---               & API Gateway centralizado (Spring Cloud Gateway). Enruta y balancea las peticiones. \\[4pt]
  uamishop-frontend      & 3000 & ---               & Interfaz de usuario est\'atica (HTML/CSS/JS) servida por Nginx. \\
\end{longtable}

\subsection{Infraestructura de Soporte}

\begin{itemize}
  \item \textbf{MySQL 8:} Cada microservicio con estado dispone de su propia instancia de
        base de datos dentro de la red Docker, garantizando aislamiento de datos.
  \item \textbf{RabbitMQ 3 (management):} Broker de mensajer\'ia configurado con un
        \textit{topic exchange} llamado \texttt{uamishop.exchange}. Accesible v\'ia
        consola de administraci\'on en el puerto 15672.
  \item \textbf{Docker Compose:} Orquestador local que levanta la totalidad de los
        contenedores con sus dependencias, redes y vol\'umenes persistentes.
\end{itemize}

%-----------------------------------------------------------------------
% 3. PATRONES IMPLEMENTADOS
%-----------------------------------------------------------------------
\section{Patrones de Dise\~no Implementados}

\subsection{Domain-Driven Design (DDD) y Bounded Contexts}

Cada microservicio representa un \textit{Bounded Context} bien definido dentro del
dominio del e-commerce:

\begin{itemize}
  \item \textbf{Cat\'alogo:} Gestiona el agregado \texttt{Producto} con sus propiedades
        (SKU, nombre, descripci\'on, precio, stock, categor\'ia). La validaci\'on del
        SKU es una regla de negocio del dominio.
  \item \textbf{Ordenes:} Gestiona el ciclo de vida de una \texttt{Orden}, desde su
        creaci\'on (\texttt{CREADA}) hasta su confirmaci\'on y pago (\texttt{PAGO\_PROCESADO}).
  \item \textbf{Ventas:} Gestiona el agregado \texttt{Carrito} con sus \texttt{ItemCarrito}.
        Coordina con el cat\'alogo para obtener datos de producto.
\end{itemize}

La arquitectura hexagonal (puertos y adaptadores) se aplica en cada servicio para aislar
la l\'ogica de dominio de los detalles de infraestructura (JPA, REST, RabbitMQ).

\subsection{Circuit Breaker con Resilience4j}

El servicio \texttt{uamishop-ventas} implementa un \textit{Circuit Breaker} sobre las
llamadas al servicio \texttt{uamishop-catalogo}. Cuando el cat\'alogo no est\'a disponible
o responde con errores, el circuit breaker se abre y la operaci\'on de agregar productos
al carrito retorna un error controlado (\texttt{HTTP 503}) en lugar de propagar el fallo.

Configuraci\'on clave en \texttt{application.yml} de ventas:
\begin{itemize}
  \item \textbf{slidingWindowSize:} 5 peticiones para calcular la tasa de fallos.
  \item \textbf{failureRateThreshold:} 50\% de fallos activa el estado OPEN.
  \item \textbf{waitDurationInOpenState:} 10 segundos antes de pasar a HALF-OPEN.
\end{itemize}

\subsection{Outbox Pattern}

El servicio \texttt{uamishop-ordenes} implementa el patr\'on \textit{Outbox} para
garantizar consistencia eventual entre la base de datos y el broker de mensajer\'ia.
Cuando se confirma o paga una orden, el evento se persiste primero en una tabla
\texttt{outbox\_events} dentro de la misma transacci\'on de base de datos y,
posteriormente, un proceso \textit{scheduler} lo publica en RabbitMQ. Esto evita
la p\'erdida de mensajes en caso de fallos parciales.

%-----------------------------------------------------------------------
% 4. API GATEWAY
%-----------------------------------------------------------------------
\section{API Gateway}

El API Gateway centraliza el acceso a todos los microservicios desde un \'unico punto
de entrada (\texttt{localhost:8090}). Las rutas est\'an configuradas en
\texttt{uamishop-gateway/src/main/resources/application.yml}.

\begin{longtable}{@{}p{5cm}p{2.5cm}p{6cm}@{}}
  \toprule
  \textbf{Patr\'on de ruta} & \textbf{Destino} & \textbf{Descripci\'on} \\
  \midrule
  \endfirsthead
  \toprule
  \textbf{Patr\'on de ruta} & \textbf{Destino} & \textbf{Descripci\'on} \\
  \midrule
  \endhead
  \bottomrule
  \endfoot
  \texttt{/api/v1/productos/**}   & catalogo:8081  & CRUD de productos del cat\'alogo. \\[4pt]
  \texttt{/api/v1/categorias/**}  & catalogo:8081  & Gesti\'on de categor\'ias de productos. \\[4pt]
  \texttt{/api/v1/carritos/**}    & ventas:8083    & Operaciones de carrito de compra. \\[4pt]
  \texttt{/api/v1/ordenes/**}     & ordenes:8082   & Creaci\'on y consulta de \'ordenes. \\
\end{longtable}

%-----------------------------------------------------------------------
% 5. RABBITMQ
%-----------------------------------------------------------------------
\section{Flujo de Eventos con RabbitMQ}

RabbitMQ se utiliza para la comunicaci\'on as\'incrona entre microservicios. Se configura
un \textit{topic exchange} llamado \texttt{uamishop.exchange}.

\begin{longtable}{@{}p{4.5cm}p{4.5cm}p{4.5cm}@{}}
  \toprule
  \textbf{Cola (Queue)} & \textbf{Routing Key} & \textbf{Descripci\'on} \\
  \midrule
  \endfirsthead
  \toprule
  \textbf{Cola (Queue)} & \textbf{Routing Key} & \textbf{Descripci\'on} \\
  \midrule
  \endhead
  \bottomrule
  \endfoot
  \texttt{ordenes.queue}         & \texttt{orden.confirmada}   & Notifica que una orden ha sido confirmada por el usuario. \\[4pt]
  \texttt{ordenes.pago.queue}    & \texttt{orden.pago}         & Notifica que el pago de una orden ha sido procesado. \\[4pt]
  \texttt{ventas.queue}          & \texttt{venta.procesada}    & Notifica que una venta ha sido completada exitosamente. \\
\end{longtable}

El patr\'on Outbox en \texttt{uamishop-ordenes} garantiza que cada evento sea publicado
exactamente una vez, incluso ante fallos transitorios en la conexi\'on con RabbitMQ.

%-----------------------------------------------------------------------
% 6. FRONTEND
%-----------------------------------------------------------------------
\section{Frontend}

El frontend de UAMIShop es una aplicaci\'on web est\'atica desarrollada con tecnolog\'ias
nativas del navegador, sin dependencias de frameworks JavaScript externos.

\begin{itemize}
  \item \textbf{HTML5 / CSS3:} Estructura y estilos de la interfaz de usuario.
  \item \textbf{JavaScript (Vanilla ES6+):} L\'ogica de interacci\'on, llamadas a la API
        REST a trav\'es del Gateway y actualizaci\'on din\'amica del DOM.
  \item \textbf{Nginx:} Servidor web que sirve los archivos est\'aticos del frontend
        en el puerto 3000 y act\'ua como proxy inverso hacia el API Gateway.
\end{itemize}

Las funcionalidades implementadas en la interfaz incluyen:

\begin{itemize}
  \item Listado y b\'usqueda de productos del cat\'alogo.
  \item A\~nadir productos al carrito de compras.
  \item Visualizaci\'on del contenido del carrito.
  \item Confirmaci\'on y pago de \'ordenes.
  \item Consulta del estado de \'ordenes existentes.
\end{itemize}

%-----------------------------------------------------------------------
% 7. DESPLIEGUE
%-----------------------------------------------------------------------
\section{Instrucciones de Despliegue}

El proyecto se despliega completamente con Docker Compose. A continuaci\'on se describen
los pasos necesarios para levantar el sistema en un entorno local.

\subsection{Requisitos Previos}

\begin{itemize}
  \item Docker Engine 24+ y Docker Compose v2+.
  \item Java 17+ y Maven 3.8+ (s\'olo para compilaci\'on sin Docker).
  \item Puertos libres: 3000, 8080, 8081, 8082, 8083, 8090, 5672, 15672, 3306.
\end{itemize}

\subsection{Comandos de Despliegue}

\begin{lstlisting}[language=bash, caption={Levantar el sistema completo}]
# Clonar el repositorio
git clone <url-repositorio> uamishop
cd uamishop

# Construir imagenes y levantar todos los contenedores
docker compose up --build -d

# Verificar que todos los contenedores esten corriendo
docker ps

# Ver logs de un servicio especifico
docker compose logs -f uamishop-ventas

# Detener el sistema
docker compose down

# Detener y eliminar volumenes (reset completo)
docker compose down -v
\end{lstlisting}

\begin{lstlisting}[language=bash, caption={Verificaci\'on r\'apida del sistema}]
# Health check del API Gateway
curl http://localhost:8090/actuator/health

# Listar productos via Gateway
curl http://localhost:8090/api/v1/productos

# Crear un carrito via Gateway
curl -X POST http://localhost:8090/api/v1/carritos \
  -H "Content-Type: application/json" \
  -d '{"clienteId": "550e8400-e29b-41d4-a716-446655440000"}'

# Acceder a la consola de administracion de RabbitMQ
# URL: http://localhost:15672
# Usuario: guest / Contrasena: guest
\end{lstlisting}

%-----------------------------------------------------------------------
% 8. CONCLUSIONES
%-----------------------------------------------------------------------
\section{Conclusiones}

El proyecto UAMIShop demuestra la viabilidad de implementar un sistema de e-commerce
completo bajo una arquitectura de microservicios, aplicando patrones modernos de dise\~no
de software en un contexto acad\'emico.

Los principales logros del proyecto son:

\begin{enumerate}
  \item \textbf{Separaci\'on de responsabilidades:} Cada microservicio tiene un dominio
        claramente delimitado, facilitando el mantenimiento y la escalabilidad independiente
        de cada componente.

  \item \textbf{Resiliencia:} La implementaci\'on del Circuit Breaker en el servicio de
        ventas garantiza que el sistema contin\'ue operativo incluso cuando el servicio de
        cat\'alogo experimenta problemas, retornando errores controlados al cliente.

  \item \textbf{Consistencia eventual:} El patr\'on Outbox en el servicio de \'ordenes
        asegura que ning\'un evento de negocio se pierda, manteniendo la coherencia entre
        la base de datos y el sistema de mensajer\'ia.

  \item \textbf{Calidad verificada:} Las pruebas de integraci\'on con Karate validan el
        comportamiento end-to-end del sistema a trav\'es del API Gateway, con una tasa de
        \'exito del 90.9\% (10 de 11 escenarios, siendo el fallo restante una condici\'on
        de carrera relacionada con el estado del circuit breaker entre pruebas consecutivas).

  \item \textbf{Despliegue reproducible:} La contenerizaci\'on completa con Docker Compose
        garantiza que el sistema puede levantarse de manera determinista en cualquier
        entorno compatible.
\end{enumerate}

Como trabajo futuro se identifican las siguientes oportunidades de mejora: implementaci\'on
de autenticaci\'on y autorizaci\'on (JWT/OAuth2), observabilidad distribuida con Zipkin o
Jaeger, escalado horizontal de microservicios cr\'iticos y migraci\'on a Kubernetes para
entornos de producci\'on.

\vfill
\begin{center}
  \small Universidad Aut\'onoma Metropolitana --- Unidad Iztapalapa --- Abril 2026
\end{center}

\end{document}
